\section{Introduzione}
\label{sec:introduzione}

\subsection{Motivazione e contesto}

L'agricoltura di precisione rappresenta oggi uno degli ambiti applicativi più
concreti dell'Internet of Things: sensori ambientali, sistemi di visione
artificiale e attuatori connessi permettono di monitorare le colture in tempo
reale, ridurre gli sprechi di acqua e agrofarmaci, e intervenire tempestivamente
alla comparsa di patologie. Come evidenziato anche nel percorso didattico del
corso (evoluzione da \emph{Agriculture 1.0}, basata su risorse umane e animali,
fino alla \emph{Agriculture 4.0}, caratterizzata dall'uso di intelligenza
artificiale e sistemi informativi per analizzare, elaborare e decidere), la
tendenza è quella di spostare sempre più capacità decisionale vicino alla fonte
del dato, riducendo la dipendenza dal cloud e i tempi di reazione.

\textbf{PomodorIA} nasce in questo contesto come progetto per l'esame di
\emph{Internet of Things}, con l'obiettivo di realizzare un caso di studio
completo e coerente di \textbf{Edge AI applicata alla Smart Agriculture}:
una serra domotica in grado di osservare le foglie di pianta di pomodoro,
diagnosticare automaticamente eventuali patologie tramite una rete neurale
convoluzionale (CNN), incrociare la diagnosi con parametri ambientali e
decidere in autonomia se attivare irrigazione, ventilazione o allarmi.

Il progetto riutilizza e porta a compimento pratico un lavoro precedente,
svolto nell'ambito del corso di \emph{Neural Networks}, in cui è stata
addestrata e validata una CNN capace di riconoscere 10 classi di stato
fogliare (9 patologie più la classe "sana") sul dataset pubblico
\emph{PlantVillage -- Tomato}. PomodorIA prende quel modello già allenato e
lo sottopone all'intero ciclo di vita richiesto da un vero deployment
\emph{Edge}: compressione (pruning e quantizzazione), integrazione in un
sistema sensori-attuatori, orchestrazione in un ciclo continuo, e
presentazione dei risultati tramite una dashboard di monitoraggio.

\subsection{Obiettivi del progetto}

Gli obiettivi che il progetto si pone, e che questa documentazione intende
argomentare in dettaglio, sono i seguenti:

\begin{enumerate}
    \item Dimostrare, con un \emph{Proof of Concept} interamente software, il
    ciclo \textbf{sense $\rightarrow$ think $\rightarrow$ act $\rightarrow$ log}
    tipico dei sistemi IoT, applicato ad uno scenario di Smart Agriculture.
    \item Portare un modello di Deep Learning già allenato (CNN) verso un
    ipotetico deployment su hardware a risorse limitate (Raspberry Pi),
    applicando e confrontando tecniche di compressione dei modelli
    (\emph{pruning} strutturato e \emph{quantizzazione}).
    \item Progettare un agente decisionale razionale, formalizzato secondo il
    modello \textbf{PEAS} (Performance, Environment, Actuators, Sensors)
    visto a lezione, capace di combinare l'informazione visiva (diagnosi
    della CNN) con l'informazione ambientale (temperatura, umidità, luce)
    secondo il paradigma della \textbf{Data Fusion}.
    \item Mappare esplicitamente l'architettura del sistema sul modello a
    sette livelli \textbf{IoT World Forum (IoTWF)}, per collegare le scelte
    implementative alla teoria del corso.
    \item Fornire un'interfaccia di monitoraggio (dashboard) che renda
    visibili e misurabili, in tempo reale, tutte le grandezze rilevanti del
    sistema: accuratezza delle diagnosi, latenza di inferenza, consumo di
    risorse, stato degli attuatori.
\end{enumerate}

\subsection{Perché una simulazione software}

Il progetto è dichiaratamente un \emph{Proof of Concept} privo di hardware
fisico: non vengono utilizzati un Raspberry Pi reale, una fotocamera reale, né
sensori ambientali reali. Questa scelta è stata presa consapevolmente per due
motivi. In primo luogo, l'obiettivo didattico del progetto è dimostrare la
comprensione dell'architettura e delle problematiche di un sistema Edge AI
IoT (gestione delle risorse limitate, compressione dei modelli, integrazione
sensori-attuatori-agente), aspetti che restano identici a prescindere dalla
presenza di hardware fisico. In secondo luogo, la simulazione consente di
riprodurre in modo controllato e ripetibile scenari altrimenti difficili da
ottenere con hardware reale (ad esempio un'epidemia di una specifica
patologia, o condizioni ambientali estreme), rendendo il sistema più adatto
alla verifica sperimentale e alla dimostrazione in sede di esame.

Il codice è comunque scritto secondo il principio del \emph{drop-in
replacement}: i moduli che simulano hardware (fotocamera, sensori, GPIO)
espongono la stessa interfaccia che avrebbero le rispettive controparti
reali, per cui il passaggio a un deployment fisico su Raspberry Pi
richiederebbe la sostituzione dei soli moduli di I/O, senza toccare la
logica applicativa (motore di inferenza, agente decisionale, orchestratore).

\subsection{Struttura del documento}

Il documento è organizzato come segue. La Sezione~\ref{sec:contesto} richiama
i concetti teorici del corso di Internet of Things utilizzati nel progetto
(IoTWF, agenti razionali, Data Fusion, Smart Agriculture). La
Sezione~\ref{sec:architettura} descrive l'architettura complessiva del
sistema e il suo mapping sui sette livelli IoTWF. La
Sezione~\ref{sec:modello-ml} descrisce il modello di Machine Learning alla
base del sistema. La Sezione~\ref{sec:moduli} entra nel dettaglio di ciascun
modulo software. La Sezione~\ref{sec:compressione} discute la pipeline di
compressione del modello per l'Edge, incluse le problematiche tecniche
incontrate e risolte durante lo sviluppo. La
Sezione~\ref{sec:risultati} presenta i risultati sperimentali. La
Sezione~\ref{sec:esecuzione} fornisce una guida pratica all'esecuzione del
progetto. Le Sezioni~\ref{sec:limiti} e~\ref{sec:conclusioni} discutono
rispettivamente i limiti noti, gli sviluppi futuri e le conclusioni.
